\textbf{Phần 1. Trắc nghiệm lựa chọn}
\setcounterpageref{ex}{0}

\begin{ex} %[4D10N59]
Phân số chỉ số phần không tô màu trong một hình chia làm 5 phần bằng nhau, đã tô màu 2 phần là
\choice
{$\dfrac{2}{5}$}
{$\dfrac{5}{2}$}
{\True $\dfrac{3}{5}$}
{$\dfrac{3}{2}$}
\loigiai{
Phần chưa tô màu: $5 - 2 = 3$ phần $\Rightarrow \dfrac{3}{5}$.\\
\textbf{Đáp án của câu hỏi là C}
}
\end{ex}

\begin{ex} %[4D10N59]
Thương của phép chia $5 : 8$ viết dưới dạng phân số là
\choice
{$\dfrac{8}{5}$}
{\True $\dfrac{5}{8}$}
{$5{,}8$}
{$8{,}5$}
\loigiai{
$5 : 8 = \dfrac{5}{8}$.\\
\textbf{Đáp án của câu hỏi là B}
}
\end{ex}

\begin{ex} %[4D10H59]
Phân số $\dfrac{16}{20}$ rút gọn về phân số tối giản là
\choice
{$\dfrac{8}{10}$}
{$\dfrac{2}{3}$}
{$\dfrac{3}{4}$}
{\True $\dfrac{4}{5}$}
\loigiai{
$\dfrac{16 : 4}{20 : 4} = \dfrac{4}{5}$.\\
\textbf{Đáp án của câu hỏi là D}
}
\end{ex}

\begin{ex} %[4D10H59]
Quy đồng mẫu số hai phân số $\dfrac{2}{5}$ và $\dfrac{3}{10}$ với mẫu số chung là 10 ta được
\choice
{$\dfrac{2}{10}$ và $\dfrac{3}{10}$}
{$\dfrac{6}{10}$ và $\dfrac{3}{10}$}
{\True $\dfrac{4}{10}$ và $\dfrac{3}{10}$}
{$\dfrac{5}{10}$ và $\dfrac{3}{10}$}
\loigiai{
$\dfrac{2 \times 2}{5 \times 2} = \dfrac{4}{10}$ và $\dfrac{3}{10}$.\\
\textbf{Đáp án của câu hỏi là C}
}
\end{ex}

\begin{ex} %[4D10H59]
Phân số nào sau đây lớn nhất?
\choice
{\True $\dfrac{6}{7}$}
{$\dfrac{3}{7}$}
{$\dfrac{5}{7}$}
{$\dfrac{2}{7}$}
\loigiai{
Cùng mẫu số: $\dfrac{6}{7}$ lớn nhất.\\
\textbf{Đáp án của câu hỏi là A}
}
\end{ex}

\begin{ex} %[4D10H59]
Phân số nào sau đây bé hơn 1?
\choice
{$\dfrac{7}{5}$}
{$\dfrac{9}{9}$}
{$\dfrac{11}{10}$}
{\True $\dfrac{5}{8}$}
\loigiai{
$\dfrac{5}{8} < 1$ (tử số nhỏ hơn mẫu số).\\
\textbf{Đáp án của câu hỏi là D}
}
\end{ex}

\begin{ex} %[4D10V59]
Tìm $x$ biết: $\dfrac{x}{15} = \dfrac{4}{5}$.
\choice
{\True $12$}
{$10$}
{$14$}
{$16$}
\loigiai{
$x = 4 \times 3 = 12$.\\
\textbf{Đáp án của câu hỏi là A}
}
\end{ex}

\begin{ex} %[4D10V59]
Sắp xếp các phân số $\dfrac{2}{3}, \dfrac{3}{4}, \dfrac{5}{6}$ theo thứ tự từ bé đến lớn:
\choice
{$\dfrac{5}{6} < \dfrac{3}{4} < \dfrac{2}{3}$}
{$\dfrac{3}{4} < \dfrac{2}{3} < \dfrac{5}{6}$}
{\True $\dfrac{2}{3} < \dfrac{3}{4} < \dfrac{5}{6}$}
{$\dfrac{2}{3} < \dfrac{5}{6} < \dfrac{3}{4}$}
\loigiai{
MSC 12: $\dfrac{8}{12} < \dfrac{9}{12} < \dfrac{10}{12} \Rightarrow \dfrac{2}{3} < \dfrac{3}{4} < \dfrac{5}{6}$.\\
\textbf{Đáp án của câu hỏi là C}
}
\end{ex}

\begin{ex} %[4D10V59]
Một lớp học có 32 học sinh, trong đó có 24 học sinh thích môn Toán. Phân số chỉ phần học sinh thích môn Toán được rút gọn tối giản là
\choice
{$\dfrac{6}{8}$}
{\True $\dfrac{3}{4}$}
{$\dfrac{12}{16}$}
{$\dfrac{2}{3}$}
\loigiai{
$\dfrac{24 : 8}{32 : 8} = \dfrac{3}{4}$.\\
\textbf{Đáp án của câu hỏi là B}
}
\end{ex}

\begin{ex} %[4D10V59]
Điền dấu thích hợp vào chỗ chấm: $\dfrac{5}{9} \dots \dfrac{5}{11}$.
\choice
{\True $>$}
{$<$}
{$=$}
{Không so sánh được}
\loigiai{
Cùng tử số: $9 < 11 \Rightarrow \dfrac{5}{9} > \dfrac{5}{11}$.\\
\textbf{Đáp án của câu hỏi là A}
}
\end{ex}

\begin{ex} %[4D10V59]
Tìm phân số bằng phân số $\dfrac{2}{3}$ có mẫu số là 24.
\choice
{$\dfrac{14}{24}$}
{\True $\dfrac{16}{24}$}
{$\dfrac{18}{24}$}
{$\dfrac{12}{24}$}
\loigiai{
$\dfrac{2 \times 8}{3 \times 8} = \dfrac{16}{24}$.\\
\textbf{Đáp án của câu hỏi là B}
}
\end{ex}

\begin{ex} %[4D10C59]
Tìm phân số tối giản $\dfrac{a}{b}$ biết nếu cộng thêm 6 vào tử số thì được phân số mới có giá trị bằng 1, và phân số $\dfrac{a}{b}$ có giá trị bằng $\dfrac{3}{5}$.
\choice
{$\dfrac{6}{10}$}
{$\dfrac{12}{20}$}
{$\dfrac{3}{5}$}
{\True $\dfrac{9}{15}$}
\loigiai{
Cộng 6 vào tử được phân số bằng 1 tức là $b - a = 6$.\\
Hiệu số phần: $5 - 3 = 2$ phần $= 6 \Rightarrow 1$ phần $= 3$.\\
Tử số $a = 3 \times 3 = 9$, Mẫu số $b = 5 \times 3 = 15$. Phân số ban đầu rút gọn tối giản là $\dfrac{3}{5}$ (hoặc phân số cần tìm là $\dfrac{9}{15}$).\\
\textbf{Đáp án của câu hỏi là D}
}
\end{ex}

\textbf{Phần 2. Trắc nghiệm đúng sai}
\setcounterpageref{ex}{0}

\begin{ex} %[4D10H59]
Xét tính đúng sai về các tính chất phân số:
\choiceTF
{\True Phân số $\dfrac{12}{18}$ rút gọn tối giản là $\dfrac{2}{3}$.}
{\True $\dfrac{4}{7} = \dfrac{8}{14} = \dfrac{12}{21}$.}
{Phân số $\dfrac{8}{9}$ lớn hơn phân số $\dfrac{9}{8}$.}
{\True $\dfrac{8}{9} < 1 < \dfrac{9}{8}$.}
\loigiai{
\begin{itemize}
    \item $12/18 = 2/3$. Đúng.
    \item Phân số bằng nhau đúng.
    \item $8/9 < 1 < 9/8 \Rightarrow 8/9 < 9/8$. Sai.
    \item So sánh qua 1 đúng.
\end{itemize}
\begin{itemize}
    \item a) \textbf{Đúng}.
    \item b) \textbf{Đúng}.
    \item c) \textbf{Sai}.
    \item d) \textbf{Đúng}.
\end{itemize}
}
\end{ex}

\begin{ex} %[4D10V59]
Cho các phân số: $\dfrac{3}{5}, \dfrac{5}{4}, \dfrac{7}{7}, \dfrac{2}{9}$. Xét tính đúng sai:
\choiceTF
{\True Phân số lớn hơn 1 là $\dfrac{5}{4}$.}
{\True Phân số bằng 1 là $\dfrac{7}{7}$.}
{Có 3 phân số bé hơn 1.}
{\True Có 2 phân số bé hơn 1 (là $\dfrac{3}{5}$ và $\dfrac{2}{9}$).}
\loigiai{
\begin{itemize}
    \item Lớn hơn 1: 5/4. Đúng.
    \item Bằng 1: 7/7. Đúng.
    \item Bé hơn 1: chỉ có 2 phân số ($3/5, 2/9$). Ý 3 Sai, Ý 4 Đúng.
\end{itemize}
\begin{itemize}
    \item a) \textbf{Đúng}.
    \item b) \textbf{Đúng}.
    \item c) \textbf{Sai}.
    \item d) \textbf{Đúng}.
\end{itemize}
}
\end{ex}

\textbf{Phần 3. Trả lời ngắn}
\setcounterpageref{ex}{0}

\begin{ex} %[4D10H59]
Rút gọn phân số $\dfrac{18}{24}$ về tối giản (dạng a/b).
\shortans{3/4}
\loigiai{
\textbf{Đáp số:} $3/4$\\
$3/4$.
}
\end{ex}

\begin{ex} %[4D10H59]
Tìm $x$ biết $\dfrac{x}{20} = \dfrac{3}{5}$.
\shortans{12}
\loigiai{
\textbf{Đáp số:} $12$\\
$x = 3 \times 4 = 12$.
}
\end{ex}

\begin{ex} %[4D10V59]
Quy đồng mẫu số hai phân số $\dfrac{3}{4}$ và $\dfrac{5}{6}$ với mẫu số chung nhỏ nhất 12. Tìm tổng hai tử số mới.
\shortans{19}
\loigiai{
\textbf{Đáp số:} $19$\\
$3 \times 3 + 5 \times 2 = 9 + 10 = 19$.
}
\end{ex}

\begin{ex} %[4D10V59]
Tìm phân số lớn nhất trong các phân số: $\dfrac{2}{3}, \dfrac{3}{5}, \dfrac{7}{15}$ (dạng a/b).
\shortans{2/3}
\loigiai{
\textbf{Đáp số:} $2/3$\\
$10/15, 9/15, 7/15 \Rightarrow 2/3$.
}
\end{ex}

\begin{ex} %[4D10V59]
Tìm phân số bé nhất trong các phân số: $\dfrac{4}{5}, \dfrac{4}{7}, \dfrac{4}{9}$ (dạng a/b).
\shortans{4/9}
\loigiai{
\textbf{Đáp số:} $4/9$\\
$4/9$.
}
\end{ex}

\begin{ex} %[4D10C59]
Tìm một phân số bằng $\dfrac{2}{5}$ có mẫu số hơn tử số 18 đơn vị. Tìm tử số của phân số đó.
\shortans{12}
\loigiai{
\textbf{Đáp số:} $12$\\
Hiệu số phần: $5 - 2 = 3$ phần $= 18 \Rightarrow 1$ phần $= 6$.\\
Tử số: $6 \times 2 = 12$.
}
\end{ex}
